\documentclass[a4paper,12pt]{article}

\usepackage[utf8]{inputenc}
\usepackage[T1]{fontenc} % per i caratteri accentati corretti in PDF
\usepackage[italian]{babel}
\usepackage{geometry}
\usepackage{setspace}
\usepackage{enumitem}
\usepackage{titlesec}
\usepackage{tocloft}
\usepackage{graphicx}
\usepackage{lmodern}
\renewcommand{\rmdefault}{lmss}

\renewcommand{\contentsname}{Indice}

\geometry{margin=2.5cm}
\setstretch{1.2}

\titleformat{\section}{\large\bfseries}{\thesection}{1em}{}
\titleformat{\subsection}{\mdseries\bfseries}{\thesubsection}{1em}{}

\begin{document}

\begin{center}
    \includegraphics[width=9.5cm]{../../../assets/logo7zpus.jpg}\\
    \Large \textbf{Verbale Interno Gruppo di Progetto}\\
    \vspace{0.5cm}
    \small\hspace{10cm} 7zpus.swe@gmail.com
\end{center}

\noindent
\textbf{Data:} 24/10/2025 \\
\textbf{Durata:} 1 ora e 20 minuti \\
\textbf{Luogo:} incontro online (Server Discord del gruppo)

\vspace{0.3cm}
\hrule
\vspace{0.5cm}

\tableofcontents

\newpage

\section{Tabella di Versionamento}
    \begin{tabular}{|c|c|c|c|}
        \hline
        \textbf{Versione} & \textbf{Data} & \textbf{Autore} & \textbf{Descrizione} \\
        \hline
        1.0 & 24/10/2025 & Fattoni Antonio & Creazione del verbale e stesura iniziale \\
        \hline
    \end{tabular}


\section{Partecipanti}
\begin{itemize}[noitemsep]
    \item Fattoni Antonio 
    \item Georgescu Diana
    \item Gingilino Aaron
    \item Laoud Zakaria
    \item Rocco Matteo Alberto
    \item Soligo Lorenzo
    \item Vigolo Davide
\end{itemize}

\section{Ordine del Giorno}
\begin{enumerate}
    \item Suddivisione tra i membri del gruppo per la redazione delle analisi dei capitolati.
    \item Discussione e definizione della stima oraria complessiva del progetto.
    \item Stima del periodo di consegna in base alle tempistiche di lavoro previste.
\end{enumerate}
\vspace{0.5cm}
\section{Svolgimento e Discussione}
Il gruppo ha discusso in merito alla suddivisione dei capitolati da analizzare. Al termine della discussione, ogni membro ha scelto e si è assunto la responsabilità di redigere l’analisi di un capitolato specifico.

Successivamente, è stata affrontata la questione della stima oraria. Ogni componente ha indicato il numero di ore che prevede di dedicare al progetto, distinguendo le ore in base ai diversi ruoli previsti. Sulla base di tali dati è stato redatto un file Excel contenente la ripartizione delle ore totali per ruolo e per membro, utile alla definizione di un preventivo complessivo realistico.

Infine, il gruppo ha discusso delle possibili tempistiche di sviluppo, tenendo conto delle disponibilità individuali e del carico di lavoro stimato. Da tale confronto è emerso un periodo di consegna indicativo, coerente con le scadenze accademiche e la pianificazione interna del gruppo.

\vspace{0.5cm}
\section{Decisioni e Conclusione}
\begin{itemize}
    \item Ogni membro del gruppo ha preso in carico un capitolato per la relativa analisi.
    \item È stato elaborato un preventivo basato sulla distribuzione delle ore per ruolo e per componente.
    \item È stato stabilito come periodo indicativo di consegna l’inizio del mese di aprile.
\end{itemize}

\subsection{Ordine del giorno Prossimo Incontro}
\begin{enumerate}
    \item Discussione sul diario di bordo
\end{enumerate}

\vfill
\begin{flushright}
    \textit{7-ZPUs}
\end{flushright}

\end{document}